\documentclass[11pt]{article}
\usepackage[utf8]{inputenc}
\usepackage{amsmath,amssymb,amsthm}
\usepackage[margin=1in]{geometry}
\usepackage{hyperref}
\newtheorem{theorem}{Theorem}
\newtheorem{proposition}[theorem]{Proposition}
\newtheorem{definition}[theorem]{Definition}
\newcommand{\Sat}{\operatorname{Sat}}
\newcommand{\Fix}{\operatorname{Fix}}
\title{EMK: Cut Records, Recognition, Native Iota and Topology\\
\large R15 corrective edition with scoped verification}
\author{Research framework: Monty Dabas}
\date{1 October 2026}
\begin{document}
\maketitle
\begin{abstract}
We correct the supplied manuscript \texttt{emk\_topology.tex} at its
cut/operator foundation. Recognition is completed against future cuts,
the Eye is correctly typed, and saturated cut stability gives a topology.
A quarter-turn operator is constructed from signed cyclic phase records,
before introducing a complex chart. Scale, seam and path memory are retained.
The original manuscript is not certified as written: its three prose axioms
do not establish its asserted four-cycle, analytic completion or physical
consequences. This edition provides written proofs and reproducible exact
checks of the stated constructions, not a proof-assistant certification.
\end{abstract}

\section{Source and proof language}
The source names the ground, a first cut $\kappa$, and a self-cut $\chi$.
The first cut initially has the ground as its domain, while later expressions
apply it to all appearances. Those extensions must be typed. We begin with
finite records of typed arrows. In the common-carrier construction below,
the admitted actions are total maps of $X$, and $r:X\to Y$ records the
self-cut output; the exact-output interpretation is $r=\chi$.
These signatures specify the interpretation being constructed. They are not
claimed to follow uniquely from the source's prose.

Logical equality, finite records, subsets, functions and induction are the
declared metamathematical language. No scalar field, metric, ordinary complex
plane, external analytic theorem, stochastic law or physical clock is an
input to the recognition/topology construction. Object-level recognition
equality is distinct from logical equality in the proof language.

\section{Typed words before a scalar chart}
\begin{theorem}[Free path construction]
Typed finite cut paths, with concatenation at matching endpoints and an
empty path at each object, have associative composition and local identities.
Assignment of generating arrows to maps with the same signatures extends
uniquely to evaluation of every path.
\end{theorem}
\begin{proof}
The empty path must map to its local identity. Evaluate a nonempty path by
evaluating its shorter prefix and then its last arrow. Induction proves
existence and uniqueness. Both parenthesizations of concatenation give the
same ordered record, proving associativity.
\end{proof}
In the common-carrier free term presentation, the four-letter word
$\chi\kappa\chi\kappa$ is not the empty word. Naming two acts does not
prove $(\chi\kappa)^2=I$. For example, on three appearance states let
$\kappa=(1,2,0)$ and $\chi=(0,1,2)$. Then $J^2(0)=2\ne0$ for
$J=\chi\kappa$. An initial arrow from the ground to the seed is separate.
This witnesses the gap in the algebraic totalization; it does not refute
another native construction whose stronger relations are actually proved.

\section{Recognition and Eye}
\begin{definition}
For the action alphabet $G$, let $G^*$ contain all finite action words,
including the empty word. Set
\[
x\approx y\quad\Longleftrightarrow\quad
r(wx)=r(wy)\quad\hbox{for every }w\in G^*.
\]
\end{definition}
\begin{theorem}[Greatest future-compatible recognition]
The relation $\approx$ is an equivalence preserved by every action. It is
the greatest action-compatible equivalence contained in $\ker(r)$.
\end{theorem}
\begin{proof}
Each equality of records is reflexive, symmetric and transitive, so their
intersection is an equivalence. Appending an action $g$ to the future test
replaces $w$ by $wg$, proving preservation. Any compatible equivalence in
$\ker(r)$ is preserved by every word by induction and therefore implies
all the displayed record equalities.
\end{proof}
The theorem covers total actions on a common carrier. Partial typed actions
also require matching future availability and source types. The finite
exhaustive certificate covers the total case. In a free term model with
literal self-cut records, the relation is equality; nontrivial collapse
requires actual source identifications or a specified readout.

Write $q:X\to Q=X/{\approx}$. The expression $q^2=q$ is not typed.
\begin{theorem}[Typed Eye and saturation]
The operator $\Sat(U)=q^{-1}(q(U))$ on subsets of $X$ is extensive,
monotone, union preserving and idempotent. Every cut descends to
$\bar g([x])=[gx]$. The zero-change seam is
\[
S_g^0=\{x:q(gx)=q(x)\}=q^{-1}(\Fix(\bar g)).
\]
\end{theorem}
\begin{proof}
Saturation includes exactly the recognition classes meeting $U$; repeating
it adds none. Inclusion and unions follow from that membership condition.
Action compatibility makes $\bar g$ well-defined. The seam identity is
then precisely the definition of a fixed point in $Q$.
\end{proof}
The original seam applies an additional $\chi$ relative to zero healing
cost. With $\chi=(0,0,1)$, $C=(2,2,2)$ and $x=0$, its seam test is true
but its healing cost is one. The original seam-zero-cost theorem is false
even in the exact-output interpretation $R=\ker\chi$.

\section{The repaired cut topology}
\begin{theorem}[Saturated cut stability]
The family
\[
\tau_{\rm cut}=\{U\subseteq X:\Sat(U)=U,\ g(U)\subseteq U\ (g\in G)\}
\]
is a topology closed under arbitrary intersections. It consists exactly of
the inverse images of forward-invariant subsets of $Q$, or equivalently
the upper sets of the preorder generated by recognition and directed cuts.
\end{theorem}
\begin{proof}
Empty and full sets satisfy the conditions. Unions preserve both conditions.
An intersection is saturated because membership in every intersectand is
constant on a class. If $x$ belongs to every intersectand, so does $gx$;
thus arbitrary intersections are stable. Saturated sets are precisely the
inverse images under $q$, and action stability descends. Closure under a
finite chain of recognition/cut steps follows by induction, proving the
preorder characterization in both directions.
\end{proof}
This is not automatically the coarsest topology making cuts continuous:
the indiscrete topology already makes all endomaps continuous. Individual
noncommuting actions need not be continuous for this invariant-set topology.
An initial topology needs specified observables and codomain topologies.

The original image condition does not imply saturation. With
$\chi=(0,0,2)$, $\kappa=(2,2,0)$ and $R=\ker\chi$, both
$\{0,2\}$ and $\{1,2\}$ satisfy it for every generated word, but their
intersection does not. The original theorem's extra direct-image
intersection-preservation premise is sufficient, but was not derived;
on all subsets it forces the quotient to be injective.

\section{Cut count and directed distance}
\begin{theorem}
Counting appended cut occurrences gives $\ell(uv)=\ell(u)+\ell(v)$.
On $Q$, the least length of a cut word from $a$ to $b$, denoted $d^+(a,b)$
and set to infinity if no such word exists, satisfies
\[
d^+(a,b)=0\iff a=b,\qquad
d^+(a,c)\le d^+(a,b)+d^+(b,c).
\]
\end{theorem}
\begin{proof}
Length additivity follows by induction on appended occurrences. Zero length
means the empty word. Concatenating finite minimizing words proves the
triangle inequality; an infinite right side is automatic. The least finite
length exists by induction on cut counts.
\end{proof}
This supplies an event index and a directed counting distance, distinct from
the original healing cost. Symmetry and physical time do not follow.

\section{Native quarter-turn from cyclic cut records}
Use the existing UGD cyclic phase alphabet of size $4m$, $m>0$. Its size is
a declared native input, not selected by the upload's A0--A2. Let $G_{4m}$
be the free signed ledger on phase marks $e_j$ with indices modulo $4m$.
Formal cancellation supplies signs; coefficients count occurrences. Set
\[
Ue_j=e_{j+1},\quad f_j=e_j-e_{j+2m},\quad
V^-=(I-U^{2m})G_{4m},\quad \iota=U^m|_{V^-},\quad Ke_j=e_{m-j}.
\]
\begin{theorem}[Operator before complex chart]
The $f_j$, $0\le j<2m$, freely generate $V^-$. On this sector,
\[
\iota^2=-I,\quad\iota^4=I,\quad K^2=I,\quad K\iota=-\iota K.
\]
\end{theorem}
\begin{proof}
Each generator image under $I-U^{2m}$ is a listed $f_j$ or its negative.
Its first-half coefficients prove independence. A half-cycle exchanges the
two terms of $f_j$, so $U^{2m}f_j=-f_j$, giving the power identities.
Reflection squares to identity and obeys $KU^m=U^{-m}K$. It preserves
$V^-$. There $U^{-m}=-U^m$, giving anticommutation.
\end{proof}
For $m=1$, $\iota f_0=f_1$ and $\iota f_1=-f_0$. Only afterward can a
two-component chart be written as $a+\iota b$. The full phase module is
larger: $U^{2m}e_0=e_{2m}\ne-e_0$. UGD scale and seam ledgers remain
separate. Extending this action by identity on a seam ledger squares to
$(-I,I)$, not global minus identity. No complex scalar replaces the full
generalized state. The construction does not identify arbitrary
$\kappa,\chi$ with these derived operators.

\section{Path memory and the involution boundary}
\begin{theorem}[Return retains its ledger]
In a four-phase directed path, the count $n=4w+r$, $0\le r<4$, gives
visible phase $r$ and loop count $w$. Concatenation is
\[
(r,w)(s,v)=((r+s)\bmod4,\ w+v+\lfloor(r+s)/4\rfloor).
\]
The empty path and the four-step path have equal visible phase and unequal
loop count. Operator identity after four steps does not imply equal histories.
\end{theorem}
\begin{proof}
Repeated removal of full four-step blocks gives the decomposition.
Adding counts gives the formula and associativity. The two displayed paths
have counts zero and four, hence loop counts zero and one.
\end{proof}
\begin{proposition}[Involution is not a helix or an Eye]
If $J^2=I$, its endpoint orbit has at most two elements. If also $J^2=J$,
then $J=I$. In any associative operator algebra $[J,J]=0$.
\end{proposition}
\begin{proof}
Even powers give $I$ and odd powers give $J$. Combining the two squared
identities gives $J=I$. The commutator is $JJ-JJ$.
\end{proof}
A two-state swap has no fixed point. The upload's chosen chart
$J(z)=i\bar z$ swaps the two signed coordinates; its fixed set is the full
diagonal, not only the positive ray. A declared four-vertex cycle has an
indiscrete forward-invariant topology, while its orbit quotient is one point.
Neither a graph cycle nor an involution supplies Maxwell equations,
thermodynamic potentials, a continuous angle or physical evolution.

\section{Evidence and source status}
The companion \texttt{EMK\_TOPOLOGY\_R15.md} provides commit-pinned native
and imported-premise citations. It credits the existing RKF future-observer
and native-operator results, UGD-1 phase/scale/seam arithmetic, and the later
native Euler/logarithm/period results in RKF and RH. The original morphic
algebra certificate explicitly has open obligations; it is not a blanket
endorsement of this upload.

\texttt{verify\_r15.py} checks the original bytes and all 48 source claim
environments, all 3678 two-action/observation systems on up to three states,
29290 saturation cases, nine boundary witnesses, seven rejected mathematical
replacements, ten unchanged-engine symbolic replays, one rejected altered
native proof, and 42 signed basis vectors across six phase alphabets.
The existing UGD-1 certificate is regenerated in memory against its old pin.
Written proofs cover the general statements; finite checks have finite scope.

\textbf{Whole-source verdict: not certified as written.}
The root obligation is a typed derivation connecting the upload's acts to
the cyclic/oriented native presentation, including all active records.
Naming the acts does not prove closure, inverses, selected phase size or
analytic completion. Those steps must be derived or precisely cited;
they cannot be replaced by a scalar chart or a new certificate label.
\end{document}
